% !TEX program = xelatex
% ============================================================================
%  东北大学机械工程与自动化学院 · 硕士学位论文或实践成果开题报告
%  LaTeX 空白模板（排版依据学院官方 Word 模板「附件6」逐项实测对齐）
%
%  一、怎么用
%    1. 全局替换 @@...@@ 占位符：封面 7 个字段 + 五个栏目的正文
%    2. 用 XeLaTeX 编译两遍（必须 XeLaTeX，不能 pdfLaTeX）：
%         xelatex 开题报告模板.tex
%       或双击 重新编译.bat
%    3. 想要网页端编辑：python latex_web_editor.py → 浏览器开 http://127.0.0.1:8766/
%    4. 上传 Overleaf：整包 zip → New Project → Upload Project → Compiler 选 XeLaTeX
%
%  二、排版规格（已与官方模板实测对齐，一般不要改）
%    页边距：上 2.50cm / 下 2.10cm / 左右 2.00cm（＝模板 70.9/59.5/56.7pt）
%    封面：学号 黑体四号；标题 黑体二号；字段名与值 宋体三号加粗（坐标逐点对齐）
%    五栏：栏标题在框外（黑体小四），提示行在框内（黑体小四）
%          框线不是表格线，而是用 eso-pic 画的「整页背景框」——所以每页都框满整页
%    工作计划：黑体五号表头，5 列
%    正文：小四宋体，行距 1.25
%    字体回退：Windows 用宋体/黑体；Overleaf 等无 Windows 字体时自动用 Fandol
%
%  三、目录约定
%    assets/template_img0.jpeg  校徽（勿删）
%    figs/                      放你自己的插图
% ============================================================================
\documentclass[12pt,a4paper]{ctexart}

\usepackage[top=2.50cm,bottom=2.10cm,left=2.00cm,right=2.00cm]{geometry}
\usepackage{array}
\usepackage{longtable}
\usepackage{graphicx}
\usepackage{makecell}
\usepackage{indentfirst}
\usepackage[hidelinks]{hyperref}
\usepackage{xeCJK}
\usepackage{eso-pic}
\usepackage{refcount}
\usepackage{calc}
\usepackage{needspace}

% 字体：本机用 Windows 宋体/黑体；无 Windows 字体（Overleaf/Linux）时自动回退
% 注意：回退分支**不依赖任何粗体字体文件**（Overleaf 的 TeX Live 里没有 FandolSong-Bold），
%      全部靠 AutoFakeBold 生成伪粗，避免 "font cannot be found" 报错
\IfFontExistsTF{SimSun}{%
  \setCJKmainfont{SimSun}[BoldFont=SimSun, BoldFeatures={AutoFakeBold=2.5}, FakeBold=2.5, ItalicFont=KaiTi]
  \setCJKsansfont{SimHei}
  \setCJKmonofont{FangSong}
  \setCJKfamilyfont{zhsongbf}[AutoFakeBold=3]{SimSun}
}{%
  \IfFontExistsTF{FandolSong-Regular.otf}{%
    \setCJKmainfont{FandolSong-Regular.otf}[AutoFakeBold=2.5, ItalicFont=FandolKai-Regular.otf]
    \setCJKsansfont{FandolHei-Regular.otf}
    \setCJKmonofont{FandolFang-Regular.otf}
    \setCJKfamilyfont{zhsongbf}[AutoFakeBold=3]{FandolSong-Regular.otf}
    % 显式绑定 ctex 的 \songti/\heiti/\kaishu，不再依赖 ctex 的平台字体探测
    \setCJKfamilyfont{zhsong}{FandolSong-Regular.otf}
    \setCJKfamilyfont{zhhei}{FandolHei-Regular.otf}
    \setCJKfamilyfont{zhkai}{FandolKai-Regular.otf}
    \renewcommand{\songti}{\CJKfamily{zhsong}}
    \renewcommand{\heiti}{\CJKfamily{zhhei}}
    \renewcommand{\kaishu}{\CJKfamily{zhkai}}
  }{%
    \setCJKmainfont{Noto Serif CJK SC}[AutoFakeBold=2.5]
    \setCJKsansfont{Noto Sans CJK SC}
    \setCJKmonofont{Noto Sans Mono CJK SC}
    \setCJKfamilyfont{zhsongbf}[AutoFakeBold=3]{Noto Serif CJK SC}
  }%
}
\setmainfont{Times New Roman}
\xeCJKDeclareCharClass{CJK}{"2460 -> "24FF}
\newcommand{\songbf}{\CJKfamily{zhsongbf}\bfseries}

\linespread{1.25}
\setlength{\parindent}{2em}
\setlength{\parskip}{2pt}
\setlength{\emergencystretch}{3em}
\renewcommand{\arraystretch}{1.15}
% ===== 页码：照官方模板（Times 10.6pt，紧贴文字区底部，封面与填表说明页不编号）=====
% 页码与框线留出自然间距（\footskip 越大离框越远；21pt ≈ LaTeX 惯用页脚位置）
\setlength{\footskip}{21pt}
% 模板：封面与填表说明页不编号，立论依据页起算为 1 → 只改显示值（物理页 -2），
% 背景框判断也走同一个 \\curpage，保证两者永远一致（切勿用 \\setcounter{page}{1}）
\newcommand{\curpage}{\number\numexpr\value{page}-1\relax}
\renewcommand{\thepage}{\curpage}
\makeatletter
\renewcommand{\ps@plain}{%
  \renewcommand{\@oddhead}{}\renewcommand{\@evenhead}{}%
  \renewcommand{\@oddfoot}{\hfil{\normalfont\fontsize{10.6pt}{13pt}\selectfont\thepage}\hfil}%
  \renewcommand{\@evenfoot}{\@oddfoot}}
\makeatother
\pagestyle{plain}

\newcolumntype{L}[1]{>{\raggedright\arraybackslash}p{#1}}
\newcommand{\BOXW}{\dimexpr\linewidth-2\tabcolsep-2\arrayrulewidth\relax}
\newcommand{\WPW}{\dimexpr\linewidth-10\tabcolsep-6\arrayrulewidth\relax}
\newcommand{\BUDW}{\dimexpr\linewidth-10\tabcolsep-6\arrayrulewidth\relax}

% ===== 封面字段行：缩进＋16pt＋值写在横线上 =====
\newlength{\labw}
\newcommand{\formline}[3]{%
  \settowidth{\labw}{\zihao{3}#2：}%
  \noindent\hspace*{#1}{\zihao{3}\songbf #2：}%
  \underline{\makebox[\dimexpr\linewidth-#1-\labw\relax][l]{\hspace{0.6em}\zihao{3}\songbf #3}}\\[20pt]}

% ===== 五栏：标题在框外；提示行在框内；框线由整页背景框绘制 =====
\newcommand{\boxtitle}[1]{\noindent{\heiti #1}\par\vspace{4pt}}
\newcommand{\boxprompt}[1]{\needspace{5\baselineskip}\noindent{\heiti #1}\par\vspace{5pt}}
\newcommand{\boxmidline}{\vspace{8pt}\noindent\makebox[0pt][l]{\hspace*{-9pt}\rule{\dimexpr\textwidth+18pt\relax}{0.4pt}}\par\vspace{10pt}}

% 整页背景框：第 #1~#2 页画框；首页顶部低 20pt（让 栏标题 留在框外）
\newcommand{\drawboxframe}[2]{%
  \ifnum\curpage<\getpagerefnumber{#1}\relax\else
  \ifnum\curpage>\getpagerefnumber{#2}\relax\else
    \ifnum\curpage=\getpagerefnumber{#1}\relax
      \AtTextLowerLeft{\hspace*{-9pt}\framebox(\dimexpr\textwidth+18pt\relax,\dimexpr\textheight-19pt\relax){}}%
    \else
      \AtTextLowerLeft{\hspace*{-9pt}\framebox(\dimexpr\textwidth+18pt\relax,\textheight){}}%
    \fi
  \fi\fi}
\AddToShipoutPictureBG{%
  \drawboxframe{bx1s}{bx1e}\drawboxframe{bx2s}{bx2e}%
  \drawboxframe{bx3s}{bx3e}\drawboxframe{bx4s}{bx4e}}
  \newenvironment{hangingrefs}
  {\par\setlength{\hangindent}{1.5em}\setlength{\parindent}{0pt}}
  {\par}

\begin{document}

% ==================== 封面 ====================
\begin{titlepage}
\thispagestyle{empty}
\vspace*{-25pt}
\begin{flushright}{\heiti\zihao{4} 学号：@@学号@@}\end{flushright}
\vspace{0.64cm}
\begin{center}
\includegraphics[width=6.84cm]{assets/template_img0.jpeg}\\[1.30cm]
{\zihao{2}\heiti 硕士学位论文或实践成果开题报告}
\end{center}
\vspace{3.85cm}

\formline{80.3pt}{论文或实践成果题目}{@@论文题目@@}
\formline{112.4pt}{姓\hspace{0.9em}名}{@@姓名@@}
\formline{117pt}{所属学科}{@@学科专业@@}
\formline{117pt}{学院（部）}{@@学院（部）@@}
\formline{117pt}{导\hspace{0.9em}师}{@@导师@@}
\formline{117pt}{副导师}{@@无则填“无”@@}

\vspace{63pt}
\begin{center}
{\zihao{3}\songbf 东北大学研究生院}\\[10pt]
{\zihao{3}\songbf 20XX 年\hspace{0.35em}XX 月\hspace{0.35em}XX日}
\end{center}
\end{titlepage}

% ==================== 填表说明 ====================
\newpage
\thispagestyle{empty}% 照官方模板：封面与填表说明页不显示页码
\begin{center}{\songti\fontsize{23.5pt}{30pt}\selectfont 填表说明}\end{center}
\vspace{0.7cm}
{\songti\fontsize{14.5pt}{22pt}\selectfont
\begin{enumerate}\setlength{\itemsep}{0.45em}
\item 本表一、二、三、四、五项在导师指导下准确、完整填写。
\item 学生在通过开题后一周内将该材料交到所在学院（部），并将电子版上传至研究生管理信息系统。
\item 学生入学后第三学期应完成论文或实践成果开题报告，按有关规定，没有完成开题报告的学生不能申请论文答辩。
\end{enumerate}
}

% ==================== 一、立论依据 ====================
\newpage
\boxtitle{一、立论依据}
\label{bx1s}
\boxprompt{课题来源、选题依据和背景情况、课题研究目的、理论意义和实际应用价值}
@@一、立论依据正文（写完后删掉下面这段提示）@@

%% ——————————————————————————————
%% 写作提示（交稿前请删除本段注释）：
%% （1）课题来源：课题来自本人岗位实践，还是导师课题？
%% （2）选题依据和背景情况：建议分"（1）政策背景""（2）现实背景"——
%%      政策背景要写到文号（如 国发〔2024〕7 号）；现实背景要有数据。
%% （3）课题研究目的：用"回答三个相互衔接的管理问题"这类句式点明目标。
%% （4）理论意义和实际应用价值：分（1）（2）（3），每条一个论点 + 一句支撑。
%% 插图（可选）：图片放 figs/ 目录，然后插入：
%%   \begin{center}\includegraphics[width=0.80\linewidth]{figs/你的图.png}\\[6pt]
%%   图1　图题（数据来源：……）\end{center}
%% ——————————————————————————————
\label{bx1e}

% ==================== 二、文献综述 ====================
\newpage
\boxtitle{二、文献综述}
\label{bx2s}
\boxprompt{国内外研究现状、发展动态；所阅文献的查阅范围及手段}
@@二、文献综述正文@@

%% ——————————————————————————————
%% 写作提示：建议分五节
%%   （一）所阅文献的查阅范围及手段（数据库、检索词、筛选标准、纳入结果）
%%   （二）设备更新决策研究现状（国外 / 国内）
%%   （三）环卫设备管理研究现状（国外 / 国内）
%%   （四）设备健康管理与预测性维护研究现状
%%   （五）总体评述与本文切入点
%% 每节末给"评述"，不要只罗列文献；引用用顺序号 [1][2]…
%% ——————————————————————————————
\label{bx2e}

% ==================== 三、研究内容 ====================
\newpage
\boxtitle{三、研究内容}
\label{bx3s}
\boxprompt{1．研究构想与思路、主要研究内容及拟解决的关键技术}
@@三、研究内容（1）研究构想与思路、主要研究内容及拟解决的关键技术@@

%% ——————————————————————————————
%% 写作提示：（1）研究构想与基本思路（2）主要研究内容——此处放论文三级目录
%%   （3）拟解决的关键技术（4）研究特色与创新点
%% 需要分隔线时用 \boxmidline（一条贯通整个框宽的横线）
%% ——————————————————————————————
\boxprompt{2．拟采取的研究方法、技术路线、实施方案及可行性分析}
@@三、研究内容（2）拟采取的研究方法、技术路线、实施方案及可行性分析@@

%% ——————————————————————————————
%% 写作提示：（1）研究方法（2）技术路线（3）实施方案（4）可行性分析
%%   可行性建议从"数据可行性 / 方法可行性 / 实施条件可行性 / 时间可行性"四条写
%% ——————————————————————————————
\label{bx3e}

% ==================== 四、研究基础 ====================
\newpage
\boxtitle{四、研究基础}
\label{bx4s}
\boxprompt{1．所需实验手段、研究条件和实验条件 }
@@四、研究基础（1）所需实验手段、研究条件和实验条件@@

%% 写作提示：写清"已有条件"与"缺口"，缺口要说明如何解决。
\boxprompt{2．所需经费，包含经费来源、开支预算（工程设备、材料须填写名称、规格、数量）}
@@四、研究基础（2）所需经费说明@@

%% ——— 经费预算表（行数可自行增删）———
\begin{center}
\begin{tabular}{|>{\centering\arraybackslash}p{0.055\BUDW}|p{0.15\BUDW}|p{0.34\BUDW}|>{\centering\arraybackslash}p{0.11\BUDW}|p{0.34\BUDW}|}
\hline
序号 & 支出科目 & 名称、规格、数量 & 金额（元） & 说明 \\
\hline
1 & @@科目@@ & @@名称、规格、数量@@ & @@金额@@ & @@说明@@ \\
\hline
2 &  &  &  &  \\
\hline
3 &  &  &  &  \\
\hline
\multicolumn{3}{|c|}{{\heiti 合计}} & @@合计@@ &  \\
\hline
\end{tabular}
\end{center}
\label{bx4e}

% ==================== 五、工作计划（这一栏本身就是表格）====================
\newpage
\noindent{\heiti 五、工作计划}\par\vspace{3pt}
\begin{longtable}{|>{\centering\arraybackslash}p{0.07\WPW}|L{0.35\WPW}|>{\centering\arraybackslash}p{0.13\WPW}|L{0.19\WPW}|L{0.26\WPW}|}
\hline
{\heiti\zihao{5} 序号} & {\heiti\zihao{5} 阶段及内容} & \makecell{{\heiti\zihao{5} 工作量估计}\\{\heiti\zihao{5} （时数）}} & {\heiti\zihao{5} 起止日期} & {\heiti\zihao{5} 阶段成果形式}\\
\hline
\endfirsthead
\hline
{\heiti\zihao{5} 序号} & {\heiti\zihao{5} 阶段及内容} & \makecell{{\heiti\zihao{5} 工作量估计}\\{\heiti\zihao{5} （时数）}} & {\heiti\zihao{5} 起止日期} & {\heiti\zihao{5} 阶段成果形式}\\
\hline
\endhead
\hline
\endfoot
\hline
\endlastfoot
1 & @@阶段及内容@@ & @@时数@@ & @@起止日期@@ & @@阶段成果形式@@ \\
\hline
2 &  &  &  &  \\
\hline
3 &  &  &  &  \\
\hline
4 &  &  &  &  \\
\hline
\hline
\multicolumn{3}{|c|}{{\heiti 合计}} & @@总时数@@ & \\
\hline
\end{longtable}

\vspace{6pt}
\noindent @@备注：如“工作量按每周投入约 20 小时、周期约 13 个月估算。”@@

\end{document}
